\chapter{Vocabulary}
\sloppy

\begin{sourcebox}
Apart from transliterated international words, the technical terms of Russian science consist almost entirely of nouns and adjectives that are formed from verbs compounded according to the rules of Chapter IV. The most important of these verbs are arranged here under approximately 75 roots.
\end{sourcebox}

\section{Plan of the chapter}

We now make a systematic study of the vocabulary of Russian mathematics by noting the extent to which the most important compound verbs (i.e. most productive of important nouns and adjectives) conform, or fail to conform, to the general framework of the preceding chapter.

In §§2, 3, 4 we begin with six roots of such importance that, apart from international words, they alone account for nearly half the vocabulary of Russian science: namely, \ru{вод-} lead, \ru{нос-} carry, \ru{ход-} go (or come), \ru{бер- (я-)} take, \ru{лаг-} lay, \ru{ста-} stand. Then (§5) come eleven roots\corrnote{原文此处写 ``ten roots''，并在 §5 开头列出五组后称 ``five others''；但实际 f 组明确列出 \ru{близкий, долгий, общий, широкий, явный, ясный} 六个形容词，因此合计为十一组。这里仅改正这一明确的计数错误。} (some of them very productive) for which, in reversal of the usual rule of deriving adjectives from verbs, the verbs are derived from a basic adjective; then (§6) twenty roots of considerable productivity, (§7) forty-one roots\corrnote{原文 §7 的标题与首句均写 ``forty roots''；逐项核对该节从 \ru{би-} 到 \ru{част-} 的根词标题，实际共有 41 组，故改正为 forty-one。} of moderate productivity, and finally (§8) a few terms that cannot be conveniently fitted into our general scheme.

For learning Russian vocabulary it is essential to keep firmly in mind the basic principle that the English word and the Russian word will usually be formed from the same Latin word, the English by transliteration and the Russian by loan-translation. In this way the task of learning several thousand Russian scientific terms is reduced to learning about 75 roots. For example, when the reader realizes that the Russian \ru{лаг-} lay, put (see §4) corresponds to the Latin \emph{pos-} (as in \emph{position}) he will automatically have learned such Russian words as \ru{изложе́ние} \emph{exposition}, \ru{положе́ние} \emph{position}, \ru{предложе́ние} \emph{proposition}, \ru{предположе́ние} \emph{supposition}, \ru{разложе́ние} \emph{decomposition}, \ru{сло́жный} \emph{composite}, etc. The present chapter will contain several hundred such examples, with a clear indication (sometimes by italicization) of both the Latin and the Russian root (in a few cases the similarity in meaning is not due to loan-translation but to historical coincidence). The reader should take active note of each of these roots whenever they occur.

\section{The three ``verbs of motion''}

The three ``verbs of motion'' \ru{води́ть} lead (Lat. \emph{duc-}), \ru{носи́ть} carry (Lat. \emph{fer-, lat-}) and \ru{ходи́ть} go, come (various Latin roots) can be compounded with each of the 19 prefixes (two or three of the 57 possibilities are rare or non-existent) to form verb-compounds that are common in Russian science, together with a constantly increasing number of derived nouns and adjectives. Under the headings a, b, c below we examine some of the mathematically important compounds of these three verbs, giving only the imperfective forms. Note that, in contrast to the usual rule (Chap. IV, §3), the lexical compounds of these verbs of motion: \ru{вводи́ть} introduce, \ru{вноси́ть} insert, \ru{входи́ть} enter, etc. are all imperfective. Their perfective partners are formed from three other verbs, which we shall not discuss here (nor will they appear in the Exercises or Readings), although one or two of their compounds occur fairly often in mathematics. In case of need the reader will be able to find them in a Russian-English dictionary.

\subsection*{a. Compounds and phrases from the root \ru{вод-} --- lead}

\begingroup
\small
\setlength{\tabcolsep}{3pt}
\renewcommand{\arraystretch}{1.08}
\begin{longtable}{@{}>{\cyrillicfont}p{27mm}p{27mm}>{\cyrillicfont}p{49mm}p{38mm}@{}}
\toprule
\normalfont Compound & Meaning & \normalfont Illustrative phrase & Meaning\\
\midrule
\endhead
вводить & introduce & вводить новую систему & introduce a new system\\
возводить & raise & возводить в степень & raise to a power\\
выводить & deduce & вывод --- ряд формул & a deduction is a sequence of formulas\\
переводить & translate & машинный перевод & machine translation\\
проводить & produce, carry out & проводить преобразование & carry out a transformation\\
сводить & reduce & аксиома сводимости & axiom of reducibility\\
\bottomrule
\end{longtable}
\endgroup

\begin{keybox}[Exercise 5.1]
Read aloud the following sentences illustrating respectively these derivatives from the root \ru{вод-} (Latin: \emph{duc-}). Note the loan-translations.
\end{keybox}

\noindent\ru{введе́ние} introduction; \ru{выводи́ть} deduce; \ru{вы́вод} deduction; \ru{приводи́ть} reduce to; \ru{неприводи́мый} irreducible; \ru{поведе́ние} conduct, behaviour; \ru{производи́ть} produce; \ru{произведе́ние} product; \ru{произво́дная} derivative; \ru{своди́ть} reduce.

\begin{enumerate}[label=\arabic*)]
\item \ru{введе́ние иррациона́льных чи́сел по Дедеки́нду}\\
      the introduction of irrational numbers according to Dedekind.
\item \ru{из фо́рмулы выво́дится но́вая фо́рмула путём подстано́вки}\\
      from a formula a new formula is deduced by means of substitution.
\item \ru{доказа́тельство теоре́мы состои́т из вы́вода фо́рмулы}\\
      the proof of a theorem consists of the deduction of a formula.
\item \ru{э́то усло́вие приво́дит к уравне́нию}\\
      this condition reduces to the equation.
\item \ru{число́ называ́ется алгебраи́ческим, е́сли оно́ явля́ется ко́рнем неприводи́мого многочле́на}\\
      a number is called algebraic if it is a root of an irreducible polynomial.
\item \ru{мы изуча́ем поведе́ние фу́нкции}\\
      we study the behavior of the function.
\item \ru{мы говори́м, что сече́ние произво́дится рациона́льным число́м}\\
      we say that the cut is produced by the rational number.
\item \ru{скаля́рным произведе́нием двух ве́кторов}\\
      by the scalar product of two vectors.
\item \ru{э́тот преде́л называ́ется произво́дной фу́нкции}\\
      this limit is called the derivative of the function.
\item \ru{на терминоло́гии $\varepsilon$--$\delta$ непреры́вность фу́нкции сво́дится к сле́дующему}\\
      in the terminology $\varepsilon$--$\delta$, continuity of a function reduces to the following.
\end{enumerate}

\subsection*{b. Compounds and phrases from the root \ru{нос-} --- carry}

\begingroup
\small
\setlength{\tabcolsep}{3pt}
\renewcommand{\arraystretch}{1.08}
\begin{longtable}{@{}>{\cyrillicfont}p{27mm}p{28mm}>{\cyrillicfont}p{45mm}p{40mm}@{}}
\toprule
\normalfont Compound & Meaning & \normalfont Illustrative phrase & Meaning\\
\midrule
\endhead
вносить & insert & вносить данные & insert data\\
выносить & carry out, export & выноска & marginal note\\
заносить & enter (in a list) & заносить законы в список & enter the laws in a list\\
наносить & carry onto, plot & наносить кривую & plot a curve\\
относить & refer, relate & относительное положение & relative position\\
переносить & transfer & перенос осей & transfer of axes\\
сносить & carry down, remove & сноска & footnote\\
\bottomrule
\end{longtable}
\endgroup

\begin{keybox}[Exercise 5.2]
Read aloud the following sentences illustrating these derivatives from the root \ru{нос-} carry (Latin: \emph{-fer, -lat-}).
\end{keybox}

\noindent\ru{носи́тель} carrier; \ru{относи́ть} relate; \ru{отноше́ние} relation; \ru{относи́тельно} relatively; \ru{относи́тельность} relativity; \ru{переноси́ть} transfer; \ru{перено́с} transference, translation; \ru{соотноше́ние} relationship.

\begin{enumerate}[label=\arabic*)]
\item \ru{замыка́ние мно́жества то́чек $x$, для кото́рых $f(x)\ne0$, называ́ется её носи́телем}\\
      the closure of the set of points $x$ for which $f(x)\ne0$ is called its carrier.
\item \ru{теоре́ма, относя́щаяся к то́чной ве́рхней грани́це}\\
      the theorem relating to the exact upper bound.
\item \ru{что́бы отноше́ние бы́ло равно́ за́данному числу́}\\
      in order that the ratio may be equal to the given number.
\item \ru{положе́ние но́вых осе́й относи́тельно ста́рой систе́мы}\\
      the position of the new axes relatively to the old system.
\item \ru{обобще́ние рима́новой геоме́трии применя́ется в о́бщей тео́рии относи́тельности}\\
      a generalization of Riemannian geometry is used in the general theory of relativity.
\item \ru{так называ́емая аксио́ма Архиме́да не перено́сится на любы́е упоря́доченные ко́льца}\\
      the so-called axiom of Archimedes does not carry over to arbitrary ordered rings.
\item \ru{преобразова́ние координа́т при паралле́льном перено́се осе́й}\\
      transformation of coordinates under parallel translation of axes.
\item \ru{фу́нкция непреры́вна, е́сли выполня́ется сле́дующее соотноше́ние}\\
      a function is continuous if the following relationship is fulfilled.
\end{enumerate}

\subsection*{c. Compounds and phrases from the root \ru{ход-} --- go, come}

\begingroup
\small
\setlength{\tabcolsep}{3pt}
\renewcommand{\arraystretch}{1.08}
\begin{longtable}{@{}>{\cyrillicfont}p{29mm}p{30mm}>{\cyrillicfont}p{46mm}p{38mm}@{}}
\toprule
\normalfont Compound & Meaning & \normalfont Illustrative phrase & Meaning\\
\midrule
\endhead
входить & enter & входные данные & input data\\
восходить & ascend & превосходить все числа класса & exceed all numbers of the class\\
выходить & turn out, issue & выходные данные & output data\\
исходить & come out, start from & исходные данные & initial data\\
находить & come upon, find & находится число & a number is found\\
обходить & go around, avoid & необходимое условие & necessary (unavoidable) condition\\
отходить & go away, digress & отходить от темы & digress from the theme\\
переходить & go across, transit & вероятность перехода & transition probability\\
подходить & go up to, approach; be suitable & подходящим образом & in a suitable manner\\
приходить(ся) & arrive; be necessary & приходится доказать & it is necessary to prove\\
проходить & go through & прямая, проходящая через точку & a line passing through a point\\
расходиться & diverge & расходящийся ряд & divergent series\\
сходиться & converge & сходящийся ряд & convergent series\\
\bottomrule
\end{longtable}
\endgroup

\begin{keybox}[Exercise 5.3]
Read aloud these sentences illustrating the following derivatives from the root \ru{ход-} go (numerous corresponding Latin roots: e.g. \emph{cess-, it-, verg-}).
\end{keybox}

\noindent\ru{входи́ть} enter; \ru{выходи́ть} issue; \ru{исходи́ть} start; \ru{исхо́дный} initial; \ru{находи́ть} find; \ru{нахожде́ние} finding; \ru{необходи́мый} necessary; \ru{переходи́ть} go over; \ru{перехо́д} transition; \ru{подходи́ть} approach; \ru{превосходи́ть} exceed; \ru{приходи́ться} be necessary; \ru{проходи́ть} go through; \ru{прохожде́ние} passage; \ru{расходи́ться} diverge; \ru{сходи́ться} converge; \ru{сходи́мость} convergence.

\begin{enumerate}[label=\arabic*)]
\item \ru{элеме́нт $-1$, входя́щий в э́ту гру́ппу, име́ет поря́док 2}\\
      the element $-1$ entering this group has order 2.
\item \ru{ве́кторы, выходя́щие из э́той то́чки}\\
      vectors issuing from this point.
\item \ru{исходя́ из не́которого значе́ния незави́симой переме́нной}\\
      starting from some value of the independent variable.
\item \ru{за исхо́дный элеме́нт построе́ния по́ля принима́ем \ldots}\\
      for the initial element of the construction of the field we take \ldots
\item \ru{нахо́дим сле́дующий результа́т}\\
      we find the following result.
\item \ru{ме́тод сумми́рования Чеза́ро применя́ется для нахожде́ния фу́нкции по её ря́ду Фурье́}\\
      the method of summation of Cesàro is used for finding a function from its Fourier series.
\item \ru{для того́, что́бы э́ти уравне́ния име́ли нетривиа́льное реше́ние, необходи́мо, что́бы детермина́нт был ра́вен нулю́}\\
      in order that these equations may have a nontrivial solution, it is necessary that the determinant be equal to zero.
\item \ru{переходя́ от поля́рных координа́т к дека́ртовым}\\
      passing from polar coordinates to Cartesian coordinates.
\item \ru{перехо́д от значе́ния $x$ к значе́нию $y$ \ldots}\\
      transition from the value $x$ to the value $y$.
\item \ru{подходя́щей дро́бью называ́ется коне́чная цепна́я дробь}\\
      by a convergent is meant a finite continued fraction.
\item \ru{нахо́дим число́, превосходя́щее да́нное число́}\\
      we find a number exceeding the given number.
\item \ru{прихо́дится по́льзоваться бо́лее сло́жным определе́нием}\\
      it is necessary to use a more complicated definition.
\item \ru{пряма́я, проходя́щая че́рез да́нную то́чку}\\
      a line passing through the given point.
\item \ru{для прохожде́ния всей ли́нии потре́буется промежу́ток вре́мени}\\
      for passage of the whole line an interval of time is required.
\item \ru{рассмо́трим сумми́рование расходя́щихся рядо́в}\\
      we shall consider summation of divergent series.
\item \ru{после́довательность схо́дится к $N$}\\
      the sequence converges to $N$.
\item \ru{э́тот ряд име́ет круг сходи́мости \ldots}\\
      this series has a circle of convergence.
\end{enumerate}

\section{The root \ru{бер- (я-)} --- take}

The important verb \ru{бра́ть} take has the perfective partner \ru{взять}. The forms of \ru{бра́ть} occur in various vowel-grades, e.g. \ru{бр-} (zero-grade, as in \ru{бра́ть}), \ru{бир-} (reduced grade, as in \ru{выбира́ть} to select), \ru{бор-} (o-grade, as in \ru{вы́бор} choice), etc. The root is the same as in the Latin \emph{transfer} (e-grade), the Greek \emph{semaphore} (signal-bearer; o-grade), and the English verb \emph{bear}.

The root of the perfective partner \ru{взять} is \ru{я-} (with perfectivizing prefix \ru{вз-}, the zero-grade of \ru{воз-}). The same root occurs in an old verb \ru{*ять} take, and again (with vowel-grade \ru{-им-}) in \ru{*имать} take, which is itself obsolete (and therefore marked here with an asterisk and left unaccented, as with other obsolete verb-forms below) but which gives rise to many compounds (\ru{понима́ть} understand, etc.). The \ru{м} in \ru{-нимать} belongs to the inherited historical alternation of this verb family; \ru{име́ть} to have is not simply a modern variant of \ru{*имать}.\corrnote{原文把 \ru{-нимать}（以及 \ru{иметь}）中的 \ru{м} 解释为原始斯拉夫语鼻化元音在俄语中“插入 m 音”去鼻化的结果，并把 \ru{иметь} 称为 \ru{*имать} 的 modern variant。这个具体历史解释不正确；此处只改正这一实质性词源错误，其余段落结构保留。}

Almost all the 19 prefixes form compounds with \ru{бра́ть} and also with \ru{-ять} (which then usually becomes \ru{-нять}; cf. p. 41), so that there are two series: \ru{вы́брать, избра́ть, \ldots} and \ru{вы́нять, поня́ть, \ldots}; in the corresponding imperfectives \ru{бра́ть} becomes \ru{-бира́ть} and \ru{-ять} becomes \ru{-нима́ть}.

\subsection*{a. Compounds and phrases illustrating the roots \ru{бер-} and \ru{я-} --- take}

\begingroup
\small
\setlength{\tabcolsep}{3pt}
\renewcommand{\arraystretch}{1.06}
\begin{longtable}{@{}>{\cyrillicfont}p{39mm}p{27mm}>{\cyrillicfont}p{42mm}p{37mm}@{}}
\toprule
\normalfont Aspect pair & Meaning & \normalfont Phrase / derivative & Meaning\\
\midrule
\endhead
выбирать $\sim$ выбрать & select & выборка & sample\\
избирать $\sim$ избрать & elect & избирательная система & electoral system\\
отбирать $\sim$ отобрать & choose & аксиома отбора & axiom of choice\\
подбирать $\sim$ подобрать & choose, pick out & подобрать число & choose a number\\
разбирать $\sim$ разобрать & take apart, analyze & разбирать вопрос & analyze a question\\
собирать $\sim$ собрать & collect & сборник & collection; title of a journal\\
внимать $\sim$ внять & listen, pay attention & принимать во внимание & take into consideration\\
занимать $\sim$ занять & occupy, busy & заниматься вопросом & busy oneself with the question\\
понимать $\sim$ понять & understand & эту теорему легко понять & this theorem is easy to understand\\
принимать $\sim$ принять & take, accept & уравнение принимает вид & the equation takes the form\\
снимать $\sim$ снять & take & снимать фотографию & take a photograph\\
\bottomrule
\end{longtable}
\endgroup

\begin{keybox}[Exercise 5.4]
Read aloud the following sentences illustrating these derivatives from the root \ru{бер- (я-)} take (Latin: \emph{capt-, cept-, fer-, lect-}).
\end{keybox}

\noindent\ru{выбира́ть} select; \ru{вы́бор} selection; \ru{набо́р} collection; \ru{подобра́ть} select from among (note \ru{подо-} for \ru{под-} before the double consonant \ru{бр-}); \ru{взять} take; \ru{име́ть} have; \ru{взаи́мно} mutually (i.e. back and forth: \ru{в} to, \ru{за} back, \ru{им-} take); \ru{объём} volume (with ``separating'' hard sign after the prefix \ru{об-}; see Chap. II, §§8--9); \ru{поня́тие} concept; \ru{понима́ть} understand, get; \ru{понима́ние} understanding, conception; \ru{принима́ть} take, accept.

\begin{enumerate}[label=\arabic*)]
\item \ru{выбира́ем отре́зок, кото́рый не соде́ржит $x$}\\
      we select a segment which does not contain $x$.
\item \ru{э́то противоре́чит вы́бору числа́ $a$}\\
      this contradicts the choice of the number $a$.
\item \ru{разли́чие проявля́ется в набо́рах аксио́м}\\
      the difference appears in the sets of axioms.
\item \ru{мо́жно подобра́ть положи́тельное число́ тако́е, что \ldots}\\
      it is possible to select a positive number such that \ldots
\item \ru{причём за $a/b$ мо́жно взять любо́е рациона́льное число́}\\
      where for $a/b$ it is possible to take an arbitrary rational number.
\item \ru{то́чка $M(a,b)$ име́ет дека́ртовы координа́ты}\\
      the point $M(a,b)$ has Cartesian coordinates.
\item \ru{ве́кторы взаи́мно перпендикуля́рны}\\
      the vectors are mutually perpendicular.
\item \ru{но при заполне́нии трёхме́рного объёма \ldots}\\
      but in the filling up of a three-dimensional volume \ldots
\item \ru{с отноше́нием поря́дка в кольце́ свя́зано поня́тие положи́тельности}\\
      with the relation of order in a ring is connected the concept of positivity.
\item \ru{мо́жно понима́ть э́ти опера́ции как опера́ции над чи́слами?}\\
      is it possible to understand these operations as operations on numbers?
\item \ru{в а́лгебре тако́е понима́ние многочле́нов неудо́бно}\\
      in algebra such a conception of polynomials is inconvenient.
\item \ru{уравне́ние принима́ет сле́дующий вид}\\
      the equation takes the following form.
\end{enumerate}

\section{The roots \ru{лаг-} --- lay and \ru{ста-} --- stand}

The three Russian roots \ru{лаг-}, \ru{ста-}, \ru{сад-} mean to put (in a lying, standing, or sitting position, respectively). In the present section we consider only: a) \ru{лаг-} and b) \ru{ста-} (often translating \ru{лаг-} by \emph{set} and \ru{ста-} by \emph{set up}), since \ru{сад-}, though very productive, is less important in mathematics (but note \ru{сосе́дний} adjacent, lit. sitting together).

\subsection*{a. Compounds and phrases illustrating the root \ru{лаг-} --- lay}

For the root \ru{лаг-} lay, which by vowel-gradation and consonant-variation also appears in the forms \ru{лож-, леж-, лег-, лог-}, we begin with the question which from now on will be important for all the verbs in the chapter: how does the aspect-pair of verbs \ru{*ложи́ть $\sim$ положи́ть} to lay (for the asterisk see §3) fit into the scheme of Chapter IV?

Here we must first remark that, although \ru{положи́ть} is very common in modern Russian, the uncompounded verb \ru{*ложи́ть} is obsolete (except colloquially), its place being taken by \ru{класть}, formerly \ru{клад-ть}, so that \ru{класть $\sim$ положи́ть} is an example of an aspect-pair formed from two different roots (\ru{клад-} and \ru{лаг-}).

This phenomenon, namely that the uncompounded verb \ru{*ложи́ть} has dropped out of the language leaving behind it numerous important compounds (\ru{вложи́ть} imbed, etc.), is quite common (cf. \ru{*имать} to take §3), other examples being provided by \ru{*станови́ть} stand §4b, \ru{*полнить} fill §5b, \ru{*врати́ть} turn §6.2, \ru{*вети́ть} speak §6.3, \ru{*воли́ть} wish §6.6, \ru{*каза́ть} show §6.12, \ru{крыть}\corrnote{原文此处以星号把 \ru{крыть} 列入已退出现代俄语的无前缀动词；但 \ru{крыть} 至今仍是标准俄语动词，故去掉星号。§6.13 对同一问题作相应改正。} cover §6.13, \ru{*образи́ть} to form §6.17.

As a natural result of the disappearance of the uncompounded \ru{*ложи́ть} (with consequent intervention of the root \ru{клад-}) most of the compounds of \ru{*ложи́ть} have two imperfectives, one formed from \ru{клад-} and the other from \ru{лаг-}; for example, both \ru{вкла́дывать} and \ru{влага́ть} are imperfectives of \ru{вложи́ть} imbed. In general, the imperfective in \ru{-кладывать} is more literal and the one in \ru{-лагать} is more metaphorical; e.g. \ru{накла́дывать $\sim$ наложи́ть} to superpose (one figure on another); \ru{налага́ть $\sim$ наложи́ть} to impose (a condition). For most of the verbs in the following list any one of the three forms could go equally well in the first column (for the last verb in the list, to add, all three are given). From the more than two hundred scientific terms derived from this one root \ru{лаг-} we can here give only a few representative examples. Let us note that the intransitive verb \ru{лежа́ть} (no perf.) to lie and some of its many compounds, e.g. \ru{подлежа́ть} (no perf.) to lie under (e.g. be subject to a condition), and \ru{принадлежа́ть} (no perf.) belong to, are also common in mathematics.

\begingroup
\small
\setlength{\tabcolsep}{3pt}
\renewcommand{\arraystretch}{1.06}
\begin{longtable}{@{}>{\cyrillicfont}p{34mm}p{28mm}>{\cyrillicfont}p{48mm}p{35mm}@{}}
\toprule
\normalfont Verb & Meaning & \normalfont Phrase / derivative & Meaning\\
\midrule
\endhead
вложить & imbed, include & отображение вложения & inclusion map\\
выкладывать & remove brackets; simplify & делая выкладки & after making the calculations / simplifications\\
докладывать & put before, report & Доклады & \emph{Doklady}, title of a journal\\
изложить & set forth, explain & элементарное изложение & an elementary exposition\\
наложить & lay on & наложимые поверхности & applicable/developable surfaces\\
отложить & lay aside, postpone & отложим рассмотрение этого вопроса & we shall postpone the discussion of this question\\
полагать & lay, set & полагая $t=\sin x$ & setting $t=\sin x$\\
предлагать & propose & предлагать задачу; предложение & propose a problem; proposition\\
предполагать & suppose, conjecture & по предположению & by supposition\\
прикладывать & apply & прикладная математика & applied mathematics\\
принадлежать & belong to & $x$ принадлежит множеству & $x$ belongs to the set\\
разложить & expand, decompose & разложение по формуле бинома & binomial expansion\\
расположить & dispose, order, place & парабола расположена вправо от оси ординат & parabola situated to the right of the axis of ordinates\\
складывать / слагать / сложить & add, put together & слагаемое; сложение натуральных чисел & addend; addition of natural numbers\\
\bottomrule
\end{longtable}
\endgroup

\begin{keybox}[Exercise 5.5]
Read aloud the following sentences illustrating these derivatives from the root \ru{лаг-} lay (English \emph{lay}; Latin \emph{pos-}, etc.).
\end{keybox}

\noindent\ru{вложи́ть} impose, imbed; \ru{вложе́ние} imbedding; \ru{изложи́ть} explain; \ru{отложи́ть} lay aside, postpone; \ru{положе́ние} position; \ru{противополо́жный} opposite; \ru{положи́тельный} positive; \ru{предположи́ть} suppose; \ru{предположе́ние} supposition; \ru{приложи́ть} apply, juxtapose; \ru{приложе́ние} application; \ru{расположи́ть} arrange; \ru{разлага́ть} expand, decompose; \ru{разложи́мый} decomposable; \ru{разложи́мость} decomposability; \ru{слага́емое} addend; \ru{сло́жный} composite, complicated; \ru{принадлежа́ть} belong; \ru{лежа́ть} lie.

\begin{enumerate}[label=\arabic*)]
\item \ru{по теоре́ме о вло́женных отре́зках}\\
      by the theorem on nested intervals.
\item \ru{вложе́ние кольца́ це́лых чи́сел в по́ле рациона́льных чи́сел обобща́ется}\\
      the embedding of the ring of integers into the field of rational numbers can be generalized.
\item \ru{мы изло́жим тео́рию иррациона́льных чи́сел}\\
      we shall explain the theory of irrational numbers.
\item \ru{ве́кторы, отло́женные из то́чки $M$}\\
      the vectors laid out from point $M$.
\item \ru{положе́ние то́чки $M$ определя́ется отре́зком $OM$}\\
      the position of point $M$ is determined by segment $OM$.
\item \ru{знак $m$ противополо́жен зна́ку $c$}\\
      the sign of $m$ is opposite to the sign of $c$.
\item \ru{кольцо́ называ́ется упоря́доченным, е́сли для его́ элеме́нтов определено́ сво́йство быть положи́тельным}\\
      a ring is said to be ordered if the property of being positive is defined for its elements.
\item \ru{предполо́жим, что да́нная пряма́я не перпендикуля́рна о́си}\\
      let us suppose that the given line is not perpendicular to the axis.
\item \ru{по предположе́нию}\\
      by supposition (hypothesis, assumption).
\item \ru{ве́ктор $b$ приложе́н к концу́ ве́ктора $a$}\\
      vector $b$ is attached to the end of vector $a$.
\item \ru{в приложе́ниях к математи́ческой фи́зике}\\
      in applications to mathematical physics.
\item \ru{располо́женное кольцо́ не име́ет дели́телей нуля́}\\
      an ordered ring does not have divisors of zero.
\item \ru{натура́льное число́ разлага́ется в произведе́ние просты́х чи́сел}\\
      a natural number is decomposed into a product of prime numbers.
\item \ru{элеме́нт, не явля́ющийся просты́м, называ́ется разложи́мым}\\
      an element not being prime is called decomposable.
\item \ru{однозна́чная разложи́мость на просты́е множи́тели}\\
      unique decomposability into prime factors.
\item \ru{чи́сла $a$ и $b$ называ́ются слага́емыми}\\
      the numbers $a$ and $b$ are called addends.
\item \ru{произво́дная сло́жной фу́нкции}\\
      the derivative of a composite function.
\item \ru{ве́рхнему кла́ссу принадлежа́т все ве́рхние грани́цы}\\
      to the upper class belong all upper bounds.
\item \ru{то́чка $A$ лежи́т на прямо́й}\\
      point $A$ lies on the line.
\end{enumerate}

\subsection*{b. The root \ru{ста-} --- stand}

The root \ru{ста-} forms a transitive verb \ru{ста́вить $\sim$ поста́вить} set, and two intransitives \ru{стать} (pf.; with the impf. \ru{станови́ться}, see below) set about doing, begin, and \ru{стоя́ть} (no pf.) stand, all three with numerous compounds.

The lexical compounds (\ru{восста́вить, вста́вить}, etc.) of \ru{ста́вить $\sim$ поста́вить} form their imperfectives with \ru{-ять} in the regular way: \ru{восстановля́ть $\sim$ восстанови́ть}, \ru{вставля́ть $\sim$ вста́вить}, etc.; see Chap. IV, §4.2.

Compounds and phrases illustrating the verb \ru{ста́вить} set:

\begingroup
\small
\setlength{\tabcolsep}{3pt}
\renewcommand{\arraystretch}{1.05}
\begin{longtable}{@{}>{\cyrillicfont}p{34mm}p{28mm}>{\cyrillicfont}p{48mm}p{35mm}@{}}
\toprule
\normalfont Verb & Meaning & \normalfont Phrase & Meaning\\
\midrule
\endhead
ставить & set & естественно ставить вопрос & it is natural to ask the question\\
восстанавливать & set up, restore & восстановить перпендикуляр & erect a perpendicular\\
вставлять & set into, insert & вставить условие & insert a condition\\
выставлять & set forth, adduce & выставить аргументы & present arguments\\
доставлять & set before, cause & доставить трудности & cause difficulties\\
заставлять & put to doing, compel & заставить нас упрощать & compel us to simplify\\
отставлять & set aside, postpone & отставить проблему & put a problem aside\\
переставлять & permute, transpose & переставить элементы матрицы & transpose elements of a matrix\\
подставлять & substitute & подставляя формулу в аксиому & substituting the formula into the axiom\\
представлять & represent & представительная выборка & representative sample\\
расставлять & arrange & расстановка точек на плоскости & arrangement of points on the plane\\
составлять & compose & детерминант, составленный из коэффициентов & determinant composed of coefficients\\
сопоставлять & associate & соответствие сопоставляет число числу & correspondence associates a number with a number\\
\bottomrule
\end{longtable}
\endgroup

The perfective verb \ru{стать} start, begin has the reflexive verb \ru{станови́ться} for its imperfective partner. The non-reflexive form \ru{*станови́ть} has dropped out of modern Russian but, like \ru{*ложи́ть} (§4a), has left important compounds behind it, e.g. \ru{установи́ть} to establish, with regularly formed imperfectives, e.g. \ru{устана́вливать}; note the regular weakening of \ru{о} to \ru{а} and the insertion of \ru{л} after \ru{в} (see Rules 1 and 2 in Chap. IV, §4).

The perfective \ru{стать} has also set up its own series of compounds, e.g. \ru{переста́ть} to cease, with regular imperfectives \ru{перестава́ть} (for the imperfectivizer \ru{-ва-} see Note 2 in Chap. IV, §6). Examples from these two series of compounds (from \ru{стать} and \ru{*станови́ть}) are:

\begin{center}
\small
\begin{tabular}{@{}>{\cyrillicfont}p{43mm}p{30mm}>{\cyrillicfont}p{67mm}@{}}
доставать $\sim$ достать & get to, obtain & достаточность ``sufficiency''\\
недоставать $\sim$ недостать & be insufficient & вскрывать недостатки ``discover shortcomings''\\
оставаться $\sim$ остаться & remain & алгоритм деления с остатком ``division with remainder''\\
переставать $\sim$ перестать & cease & не переставая ``without stopping''\\
восстанавливать $\sim$ восстановить & restore & восстанавливающее устройство ``restoring component\\
останавливаться $\sim$ остановиться & stop & останавливаемся здесь ``we stop here''\\
устанавливать $\sim$ установить & establish & аналогично устанавливается, что \ldots\\
\end{tabular}
\end{center}

The imperfective \ru{стоя́ть} ``stand'' also appears in \ru{состоя́ть} ``consist'', \ru{постоя́нный} ``constant'', \ru{расстоя́ние} ``distance'', \ru{усто́йчивость} ``stability'', etc.

\begin{keybox}[Exercise 5.6]
Read aloud the following sentences illustrating these derivatives from the root \ru{ста-} stand, set, put (Latin \emph{sta-, sti-, pos-}).
\end{keybox}

\noindent\ru{подставля́ть} substitute; \ru{подстано́вка} substitution; \ru{представля́ть} represent; \ru{представле́ние} representation; \ru{составля́ть} compose, constitute; \ru{сопоставля́ть} associate; \ru{составно́й} composite; \ru{состоя́ть} consist of; \ru{доста́точный} sufficient; \ru{оста́льной} remaining; \ru{оста́ться} remain; \ru{оста́ток} remainder; \ru{остава́ться} remain; \ru{расстоя́ние} distance; \ru{устана́вливать} establish.

\begin{enumerate}[label=\arabic*)]
\item \ru{подставля́я фо́рмулу $A$ в аксио́му 10}\\
      substituting formula $A$ into axiom 10.
\item \ru{приме́ним пра́вило подстано́вки к аксио́мам 3 и 4}\\
      we apply the rule of substitution to axioms 3 and 4.
\item \ru{уравне́ние представля́ет собо́й пряму́ю}\\
      the equation represents a straight line.
\item \ru{ка́ждый элеме́нт кольца́ определя́ется его́ представле́нием}\\
      each element of the ring is determined by its representation.
\item \ru{детермина́нт, соста́вленный из коэффицие́нтов уравне́ний}\\
      the determinant composed of the coefficients of the equations.
\item \ru{сложе́ние сопоставля́ет натура́льным чи́слам $a$ и $b$ натура́льное число́ $a+b$}\\
      addition associates with the natural numbers $a$ and $b$ the natural number $a+b$.
\item \ru{число́ 1 не счита́ется ни просты́м, ни составны́м}\\
      the number 1 is considered neither prime nor composite.
\item \ru{мно́жество, состоя́щее из одного́ числа́}\\
      a set consisting of one number.
\item \ru{необходи́мое и доста́точное усло́вие}\\
      a necessary and sufficient condition.
\item \ru{ни́жний класс соде́ржит все оста́льные рациона́льные чи́сла}\\
      the lower class contains all the remaining rational numbers.
\item \ru{гла́вные криви́зны меня́ют свои́ значе́ния, но их произведе́ние остаётся постоя́нным}\\
      the principal curvatures change their values, but their product remains constant.
\item \ru{алгори́тм деле́ния с оста́тком}\\
      the algorithm of division with a remainder.
\item \ru{для ко́рня, оста́вшегося по́сле деле́ния}\\
      for the root remaining after the division.
\item \ru{расстоя́ние ме́жду то́чками}\\
      the distance between points.
\item \ru{мо́жно установи́ть взаи́мно однозна́чное соотве́тствие}\\
      it is possible to establish a one-to-one correspondence.
\end{enumerate}

\section{Verbs formed from adjectives}

In this section we discuss eleven roots: a) \ru{рав-} equal, b) \ru{полн-} full, c) \ru{прав-} right, d) \ru{мног-} many, e) \ru{один-} one, and f) six others, all of them illustrating the formation of verbs from an adjective.\corrnote{原文写 ``ten roots''，并称 f) 为 ``five others''；但 f) 实际列出六个形容词，故合计为十一组。}

\subsection*{a. \ru{рав-} --- equal}

From the adjective \ru{ра́вный} equal comes the verb \ru{равня́ть} to make level, to set equal (with the reflexive \ru{равня́ться} to be equal to, and with numerous compounds); the adjective \ru{ра́вный} in turn is formed from an earlier noun \ru{рав-} lost in Russian but preserved in many other languages, where it means a level field (cognate with the Latin \emph{rural}).

The root \ru{рав-} gives rise to an impressive list of derived nouns and adjectives, some of which are: \ru{равне́ние} equalization, \ru{раве́нство} equality, \ru{равни́на} plain, \ru{равнове́сие} equilibrium (\ru{вес} weight), \ru{равноде́йствующая} resultant (i.e. equal acting), \ru{равноде́нствие} equinox (lit. equal day), \ru{равноме́рный} uniform, \ru{равнопра́вие} equal rights, \ru{равноси́льный} equivalent, \ru{равносторо́нний} equilateral, \ru{равносходи́мость} equiconvergence, \ru{равноуго́льный} equiangular, \ru{равноудалённый} equidistant, etc., and to many compound verbs and their derivatives; e.g., from the uncompounded \ru{равня́ть} make equal come such words as \ru{прира́внивать $\sim$ приравня́ть} to equate; \ru{сра́внивать $\sim$ сравня́ть} (or \ru{сравни́ть}) to compare; \ru{сравне́ние} congruence, \ru{сравни́тельно} comparatively; \ru{ура́внивать $\sim$ уравня́ть} to equalize; \ru{уравне́ние} equation, \ru{уравнове́шивание} equilibration (\ru{веш-} from \ru{вес} weight with consonant-variation); \ru{уравнове́шивать $\sim$ уравнове́сить} to counterbalance, etc.; and also (with vowel-gradation) \ru{рове́сник} contemporary (i.e. person of the same age), \ru{ро́вно} precisely, \ru{ро́вный} flat, \ru{ро́вность} level surface, \ru{ровня́} one's equal, \ru{ровня́ть $\sim$ сровня́ть} align, \ru{вро́вень} (also \ru{наравне́}) on a level with, \ru{у́ровень} sea level, and the perfectives of several verbs meaning ``to make straight or level'' in a literal sense: \ru{выра́внивать $\sim$ вы́ровнять} to straighten out (e.g. a road), \ru{зара́внивать $\sim$ заровня́ть} to level up (e.g. a hole in the road), \ru{разра́внивать $\sim$ разровня́ть} to level off (e.g. a football field, both lengthwise and widthwise; note how \ru{раз-} usually implies ``in various directions''), \ru{сра́внивать $\sim$ сровня́ть} to make level (e.g. a road); \ru{ура́внивать $\sim$ уро́внять} to make smooth. Finally, some of these words, like many others in Russian, will admit the negative prefix \ru{не-}; thus \ru{нера́вный} unequal; \ru{нера́венство} inequality, etc.

\begin{keybox}[Exercise 5.7]
Read aloud the following sentences illustrating these derivatives from the root \ru{рав-} level, equal (Lat. \emph{equ-, par-, uni-}).
\end{keybox}

\noindent\ru{ра́вный} equal; \ru{равня́ться} be equal to; \ru{раве́нство} equality; \ru{нера́венство} inequality; \ru{равноме́рно} uniformly; \ru{сравни́ть} compare; \ru{сравни́тельно} comparatively; \ru{уравне́ние} equation.

\begin{enumerate}[label=\arabic*)]
\item \ru{е́сли оди́н из сомножи́телей ра́вен нулю́, то и произведе́ние равно́ нулю́}\\
      if one of the factors is equal to zero, then the product also is equal to zero.
\item \ru{множи́тели $\lambda(t)$ выбира́ются так, что́бы при не́котором измене́нии пара́метра $t$ равня́лись 1}\\
      the factors $\lambda(t)$ are chosen so that for some variation of the parameter $t$ they become equal to 1.
\item \ru{вычита́я раве́нство (2) из раве́нства (1)}\\
      subtracting equality (2) from equality (1).
\item \ru{все чи́сла мно́жества бу́дут удовлетворя́ть нера́венству}\\
      all numbers of the set will satisfy the inequality.
\item \ru{легко́ ви́деть, что $f_n(x)\to0$ равноме́рно}\\
      it is easy to see that $f_n(x)\to0$ uniformly.
\item \ru{сравни́м пло́щадь $\Delta P$ с площадя́ми прямоуго́льников}\\
      let us compare the area $\Delta P$ with the areas of the rectangles.
\item \ru{дробь со сравни́тельно небольши́м знамена́телем}\\
      a fraction with a comparatively small denominator.
\item \ru{уравне́ние прямо́й име́ет вид \ldots}\\
      the equation of a straight line has the form \ldots
\end{enumerate}

\subsection*{b. \ru{полн-} --- full}

From the adjective \ru{по́лный} full comes the obsolete \ru{*полнить} to fill, which has left behind it many important compounds: e.g. \ru{выполня́ть $\sim$ вы́полнить} to complete, perform (e.g. an operation) or fulfil (e.g. a condition), \ru{дополня́ть $\sim$ допо́лнить} to supplement, complement, etc.

\begin{keybox}[Exercise 5.8]
Read aloud the following sentences illustrating these derivatives from the root \ru{полн-} full (English \emph{full}, Latin \emph{plen-}; the \emph{n} in Latin and Russian is an adjectival suffix).
\end{keybox}

\noindent\ru{полнота́} completeness; \ru{вы́полнить} fulfil; \ru{выполне́ние} fulfilment; \ru{выполни́мость} fulfillability/performability; \ru{дополни́тельный} supplementary; \ru{заполне́ние} filling; \ru{вполне́} fully.

\begin{enumerate}[label=\arabic*)]
\item \ru{э́та аксио́ма называ́ется аксио́мой полноты́}\\
      this axiom is called the axiom of completeness.
\item \ru{в расшире́нии $B$ выполни́ма опера́ция, кото́рая в $A$ была́ невыполни́ма}\\
      in the extension $B$ there is a fulfillable operation which in $A$ was not fulfillable.
\item \ru{объясня́ем на приме́рах тре́бование выполни́мости}\\
      we clarify by examples the requirement of fulfillability.
\item \ru{при дополни́тельном усло́вии}\\
      under the supplementary condition.
\item \ru{при заполне́нии трёхме́рного объёма}\\
      for the filling of a three-dimensional volume.
\item \ru{ста́вится в соотве́тствие вполне́ определённый элеме́нт}\\
      a completely defined element is put into correspondence.
\end{enumerate}

\subsection*{c. \ru{прав-} --- right}

Similarly, from the adjective \ru{пра́вый} right come verbs like \ru{выправля́ть $\sim$ вы́править} to adjust, \ru{направля́ть $\sim$ напра́вить} to direct (e.g. \ru{направле́нное мно́жество} a directed set); \ru{управля́ть $\sim$ упра́вить} to control (e.g. \ru{оптима́льное управле́ние} optimal control), etc.

\begin{keybox}[Exercise 5.9]
Read aloud the following sentences illustrating these derivatives from the root \ru{прав-} right, correct (Latin \emph{rect-}).
\end{keybox}

\noindent\ru{пра́вило} rule; \ru{впра́во} to the right; \ru{направля́ть $\sim$ напра́вить} direct; \ru{направле́ние} direction; \ru{справедли́вый} correct; \ru{справедли́вость} correctness.

\begin{enumerate}[label=\arabic*)]
\item \ru{е́сли по не́которому пра́вилу ка́ждому значе́нию $x$ ста́вится в соотве́тствие значе́ние $y$}\\
      if according to some rule to each value of $x$ there is put in correspondence a value of $y$.
\item \ru{пара́бола располо́жена впра́во от о́си ордина́т}\\
      the parabola is situated to the right of the axis of ordinates.
\item \ru{ве́ктором называ́ется направле́нный отре́зок}\\
      by a vector is meant a directed segment.
\item \ru{уравне́ние прямо́й, проходя́щей че́рез да́нную то́чку в за́данном направле́нии}\\
      the equation of a line passing through a given point in an assigned direction.
\item \ru{раве́нство справедли́во}\\
      the equality is correct.
\item \ru{легко́ установи́ть справедли́вость теоре́мы}\\
      it is easy to establish the correctness of the theorem.
\end{enumerate}

\subsection*{d. \ru{мног-} --- many}

From \ru{мно́го} many come the aspect-pair \ru{мно́жить $\sim$ умно́жить} to multiply (more commonly, \ru{умножа́ть $\sim$ умно́жить}) and lexical compounds like \ru{перемножа́ть $\sim$ перемно́жить} to multiply out, \ru{размножа́ть $\sim$ размно́жить} to reproduce, multiply, and also many nouns and adjectives like \ru{сомно́житель} factor (or simply \ru{мно́житель}), \ru{мно́жество} set (lit. many-ness), \ru{многообра́зие} manifold, \ru{умноже́ние} multiplication, etc.

\begin{keybox}[Exercise 5.10]
Read aloud the following sentences illustrating these derivatives from the root \ru{мног-} many (Latin \emph{mult-}; Greek \emph{poly-}).
\end{keybox}

\noindent\ru{мно́жество} set and \ru{подмно́жество} subset; \ru{мно́житель} multiplier, factor; \ru{сомно́житель} factor/cofactor; \ru{многообра́зие} manifold; \ru{многочле́н} polynomial; \ru{многозна́чность} many-valuedness; \ru{умно́жить} multiply.

\begin{enumerate}[label=\arabic*)]
\item \ru{отображе́ние мно́жества на подмно́жество}\\
      a mapping of a set onto a subset.
\item \ru{однозна́чная разложи́мость на просты́е множи́тели}\\
      unique decomposability into prime factors.
\item \ru{е́сли оди́н из сомножи́телей ра́вен нулю́}\\
      if one of the factors is equal to zero.
\item \ru{фу́нкция, определённая на многообра́зии}\\
      a function defined on a manifold.
\item \ru{многочле́ны нахо́дятся в о́бщей тео́рии ко́лец и поле́й}\\
      polynomials occur in the general theory of rings and fields.
\item \ru{при изуче́нии многозна́чности аналити́ческих фу́нкций}\\
      in the study of the many-valuedness of analytic functions.
\item \ru{ка́ждый член ря́да умножа́ется на не́который множи́тель}\\
      each term of the series is multiplied by some factor.
\end{enumerate}

\subsection*{e. \ru{один-} --- one}

Similarly, from \ru{оди́н} one come verbs like \ru{соединя́ть $\sim$ соедини́ть} to combine (note the vowel-gradation from \ru{один} to \ru{един-}), \ru{присоединя́ть $\sim$ присоедини́ть} to adjoin, nouns like \ru{едини́ца} unity, \ru{объедине́ние} union, join, and adjectives like \ru{еди́нственный} unique.

\begin{keybox}[Exercise 5.11]
Read aloud the following sentences illustrating these derivatives from the root \ru{один-} one (Latin \emph{uni-}).
\end{keybox}

\noindent\ru{едини́ца} unity; \ru{едини́чный} unit; \ru{еди́нственный} unique; \ru{еди́нственность} uniqueness; \ru{объединя́ть} unite; \ru{объедине́ние} union; \ru{одина́ковый} identical/equal; \ru{однозна́чный} one-valued; \ru{одноро́дный} uniform/homogeneous (\ru{неодноро́дный} inhomogeneous); \ru{присоедини́ть} adjoin; \ru{присоедине́ние} adjunction; \ru{соединя́ть} unite/connect.

\begin{enumerate}[label=\arabic*)]
\item \ru{где $a$ --- положи́тельное число́, отли́чное от едини́цы}\\
      where $a$ is a positive number different from unity.
\item \ru{в ка́ждой то́чке определя́ется едини́чный каса́тельный ве́ктор}\\
      at each point there is defined a unit tangent vector.
\item \ru{пересече́ние пара́болы с о́сью симме́трии име́ется в еди́нственной то́чке}\\
      the intersection of the parabola with the axis of symmetry occurs at a unique point.
\item \ru{чем дока́зана еди́нственность нуля́}\\
      by which is proved the uniqueness of the zero.
\item \ru{существу́ет те́рмин, объединя́ющий ма́ксимум и ми́нимум --- экстре́мум}\\
      there exists a term uniting maximum and minimum---extremum.
\item \ru{объедине́нием мно́жеств $A$ и $B$ называ́ется мно́жество элеме́нтов, принадлежа́щих ли́бо $A$, ли́бо $B$}\\
      by the union of the sets $A$ and $B$ is meant the set of elements belonging either to $A$ or to $B$.
\item \ru{ве́кторы называ́ются ра́вными, е́сли они́ име́ют одина́ковые дли́ны и одина́ковые направле́ния}\\
      vectors are called equal if they have the same lengths and directions.
\item \ru{взаи́мно однозна́чным соотве́тствием называ́ется \ldots}\\
      by a one-to-one correspondence is meant \ldots
\item \ru{дифференциа́л представля́ет лине́йную, одноро́дную фу́нкцию от прираще́ния $\Delta x$}\\
      a differential represents a linear homogeneous function of the increment $\Delta x$.
\item \ru{е́сли к э́тим откры́тым промежу́ткам присоедини́ть оди́н промежу́ток}\\
      if to these open intervals we adjoin one interval.
\item \ru{присоедине́ние к постро́енному выраже́нию зна́ка $\neg$ (не)}\\
      the adjunction to the constructed expression of the symbol $\neg$ (not).
\item \ru{отре́зок, соединя́ющий верши́ны э́ллипса}\\
      the segment joining the vertices of the ellipse.
\end{enumerate}

\subsection*{f. Six further adjectives}

Finally, let us note the six adjectives\corrnote{原文此句写 ``five adjectives，但紧接着实际列出六个：\ru{близкий, долгий, общий, широкий, явный, ясный}，故改为 six。} \ru{бли́зкий} near, \ru{до́лгий} long, \ru{о́бщий} general, \ru{широ́кий} wide, \ru{я́вный} evident, \ru{я́сный} clear (\ru{нея́сный} unclear); with the derived verbs and nouns: \ru{бли́зость} nearness; \ru{приближа́ться $\sim$ прибли́зиться} (also \ru{бли́зиться $\sim$ прибли́зиться}) to approach; \ru{приближе́ние} approximation; \ru{продолжа́ть $\sim$ продо́лжить} to lengthen, extend; \ru{аналити́ческое продолже́ние} analytic continuation; \ru{обобща́ть $\sim$ обобщи́ть} to generalize; \ru{обобще́ние} generalization; \ru{расширя́ть $\sim$ расши́рить} to widen, extend; \ru{расшире́ние по́ля} extension of a field; \ru{явля́ться $\sim$ яви́ться} show oneself as, be; \ru{выясня́ть $\sim$ вы́яснить} ascertain, clear up, etc.

\begin{keybox}[Exercise 5.12]
Read aloud the following sentences illustrating these derivatives from the roots \ru{близ-} near (Lat. \emph{proxim-}), \ru{долг-} long, \ru{шир-} wide, \ru{яв-} evident, \ru{яс-} clear.
\end{keybox}

\noindent\ru{приближа́ться} approach/approximate; \ru{приближе́ние} approximation; \ru{продолже́ние} prolongation/extension/continuation; \ru{расши́рить} widen/extend; \ru{явля́ться} be/show oneself as; \ru{выявля́ть} reveal; \ru{проявля́ть} display; \ru{нея́вный} implicit; \ru{я́сный} clear; \ru{объясня́ть} clarify.

\begin{enumerate}[label=\arabic*)]
\item \ru{гипе́рбола приближа́ется к асимпто́там}\\
      the hyperbola approaches the asymptotes.
\item \ru{приближе́ние де́йствительного числа́ рациона́льными дробя́ми}\\
      approximation of a real number by rational fractions.
\item \ru{фу́нкция $f(z)$ име́ет не бо́льше одного́ аналити́ческого продолже́ния}\\
      the function $f(z)$ has not more than one analytic continuation.
\item \ru{расши́рим да́нное по́ле до по́ля, где многочле́н име́ет ко́рень}\\
      we shall extend the given field to a field in which the polynomial has a root.
\item \ru{приме́ром алгебраи́ческой опера́ции явля́ется умноже́ние подстано́вок}\\
      an example of an algebraic operation is multiplication of substitutions.
\item \ru{ско́бки выявля́ют строе́ние фо́рмул}\\
      brackets reveal the structure of formulas.
\item \ru{разли́чие проявля́ется в набо́рах аксио́м}\\
      the difference displays itself in the choices of the axioms.
\item \ru{нея́вная фу́нкция}\\
      implicit function.
\item \ru{тепе́рь я́сно, как постро́ить кольцо́ многочле́нов}\\
      now it is clear how to construct the ring of polynomials.
\item \ru{объясня́ем э́то тре́бование}\\
      let us clarify this requirement.
\end{enumerate}


\section{Twenty roots of considerable productivity}

Under this heading we shall examine the following twenty roots, arranged here in alphabetic order, although some of them are noticeably more productive than others:

\noindent 1) \ru{вер-} believe, 2) \ru{верт-} turn, 3) \ru{вет-} speak, 4) \ru{вид-} see, 5) \ru{вис-} hang, 6) \ru{вол-} wish, 7) \ru{да-} give, 8) \ru{дел-} divide, 9) \ru{зна-} know, 10) \ru{зыв-} call, 11) \ru{иск-} seek, 12) \ru{каз-} show, 13) \ru{кры-} cover, 14) \ru{мен-} change, 15) \ru{мер-} measure, 16) \ru{мог-} be able, 17) \ru{рез-} cut, 18) \ru{сек-} cut, 19) \ru{след-} follow, 20) \ru{чит-} consider.

\subsection*{1. \ru{вер-} --- believe}

The verbs derived from \ru{вер-} believe (Lat. \emph{cred-, ver-}) fit precisely into the scheme of Chapter IV. For example: \ru{нельзя́ ве́рить} it is impossible to believe; \ru{доверя́ть $\sim$ дове́рить} confide, with \ru{довери́тельный преде́л} confidence limit; \ru{проверя́ть $\sim$ прове́рить} verify, with \ru{прове́рить вычисле́ния} verify the computations.

Sentences illustrating a) \ru{ве́рный} credible, valid (\ru{неве́рный} untrue), b) \ru{вероя́тный} probable, c) \ru{прове́рить} verify:

\begin{enumerate}[label=\alph*)]
\item \ru{э́то пра́вило ве́рно и для $n$ слага́емых} --- this rule is valid also for $n$ addends;
\item \ru{не́которые матема́тики счита́ют бо́лее вероя́тным} --- some mathematicians consider it more probable;
\item \ru{легко́ прове́рить, что э́та тополо́гия удовлетворя́ет аксио́мам} --- it is easy to verify that this topology satisfies the axioms.
\end{enumerate}

\subsection*{2. \ru{верт-} --- turn}

\ru{верт-} turn; cf. Latin \emph{revert, convert, vortex}, etc. The simple verb \ru{*врати́ть} turn is obsolete (cf. p. 78) but has given rise to regularly formed compounds; e.g. \ru{обраща́ть $\sim$ обрати́ть} turn back, return; invert, convert (note that \ru{обвр-} is simplified to \ru{обр-}); \ru{обра́тный} inverse; \ru{превраща́ть $\sim$ преврати́ть} convert, transform. Compare \ru{возвра́тное уравне́ние} reciprocal equation, \ru{обрати́мая ма́трица} invertible matrix, \ru{обра́тная тригонометри́ческая фу́нкция} inverse trigonometric function, and \ru{превраще́ние Фурье́} Fourier transform.

Sentences illustrating a) \ru{вре́мя} time, orig. \ru{врет-мя}, neuter (i.e. that which is measured by returns of an indicator to the same position), with the irregular declension \ru{вре́мя, вре́мени, вре́мени, вре́мя, вре́менем, о вре́мени}; b) \ru{враще́ние} rotation; c) \ru{обраща́ться} convert itself into; d) \ru{обра́тно} conversely; e) \ru{обрати́мость} invertibility; f) \ru{поворо́т} rotation:

\begin{enumerate}[label=\alph*)]
\item \ru{рассчита́ем вре́мя, необходи́мое для прохожде́ния вдоль ли́нии} --- let us calculate the time necessary for passage along the line;
\item \ru{ско́рость враще́ния ве́ктора} --- the speed of rotation of the vector;
\item \ru{определе́ние изоморфи́зма обраща́ется в определе́ние эквивале́нтности} --- the definition of isomorphism converts into the definition of equivalence;
\item \ru{и обра́тно} --- and conversely;
\item \ru{сложе́ние име́ет сво́йство обрати́мости} --- addition has the property of invertibility;
\item \ru{преобразова́ние координа́т при поворо́те осе́й} --- transformation of coordinates under rotation of the axes.
\end{enumerate}

\subsection*{3. \ru{вет-} --- speech}

The meaning of the old Russian noun \ru{*вет} is consultation (cf. \ru{сове́т} soviet, workers' council). The corresponding verb \ru{*вети́ть} speak has dropped out of the language (see p. 78), but there remain a few compounds, which are fairly regular from the point of view of Chapter IV; e.g. \ru{отвеча́ть $\sim$ отве́тить} answer, correspond; \ru{соответствова́ть} (no pf.) correspond.

Sentences illustrating a) \ru{отвеча́ть} answer, correspond, b) \ru{соответствова́ть} correspond, c) \ru{соотве́тствие} correspondence, d) \ru{соотве́тственно} correspondingly, respectively:

\begin{enumerate}[label=\alph*)]
\item \ru{ма́лому прираще́нию аргуме́нта отвеча́ет ма́лое прираще́ние фу́нкции} --- to a small increment of the argument corresponds a small increment of the function;
\item \ru{расстоя́ние ме́жду ордина́тами асимпто́т и соответству́ющими ордина́тами криво́й} --- the distance between the ordinates of the asymptotes and the corresponding ordinates of the curve;
\item \ru{ка́ждому значе́нию $x$ ста́вится в соотве́тствие одно́ значе́ние $y$} --- to each value of $x$ one value of $y$ is put in correspondence;
\item \ru{обозна́чим мо́дули ве́кторов $a$ и $b$ соотве́тственно че́рез $|a|$ и $|b|$} --- we denote the moduli of the vectors $a$ and $b$ respectively by $|a|$ and $|b|$.
\end{enumerate}

\subsection*{4. \ru{вид-} --- see}

The uncompounded verb is \ru{ви́деть $\sim$ уви́деть} see (Lat. \emph{vid-, vis-}). The only derived words common in mathematics are the noun \ru{вид} aspect, form and its compounds \ru{видоизменя́ть} modify (lit. change the form of) and \ru{очеви́дный} obvious (lit. to the eye visible).

Sentences illustrating a) \ru{ви́деть} see, b) \ru{вид} form, c) \ru{ввиду́} in view of:

\begin{enumerate}[label=\alph*)]
\item \ru{мы ви́дим, что пряма́я определя́ется уравне́нием пе́рвой сте́пени} --- we see that a straight line is defined by an equation of the first degree;
\item \ru{уравне́ние прямо́й име́ет вид \ldots} --- the equation of a straight line has the form \ldots;
\item \ru{ввиду́ того́, что уравне́ние соде́ржит то́лько квадра́ты переме́нных} --- in view of the fact that the equation contains only squares of the variables.
\end{enumerate}

\subsection*{5. \ru{вис-} --- hang}

The uncompounded verb is \ru{висе́ть} (no pf.) hang, depend (Latin \emph{pend-}), from the noun \ru{вес} weight. Compounds are \ru{зави́сеть} (no pf.) be dependent, as in \ru{зави́симая и незави́симая переме́нные} the dependent and independent variables, \ru{зави́симость} dependence, \ru{величина́ меня́ется в зави́симости от $x$}; \ru{ве́сить} (no pf.) weigh, as in \ru{весова́я фу́нкция} weight function; \ru{взве́шивать $\sim$ взве́сить} to weight, as in \ru{взве́шенный отбо́р} weighted sample. Note the zero vowel-grade in \ru{вз-} for \ru{воз-} and the consonant-variation in \ru{веш-} for \ru{вес-}.

\subsection*{6. \ru{вол-} --- wish}

Again the uncompounded verb \ru{*воли́ть} wish, cognate with the Latin \emph{volition}, etc., is obsolete. Common in mathematics are the adjective \ru{произво́льный} arbitrary, the adverbs \ru{дово́льно} enough (i.e. up to one's wish) and \ru{про-из-во́ль-но} at will, arbitrarily, and the compound verbs \ru{по-з-вол-я́ть} (\ru{з} for \ru{из}) allow, enable (e.g. \ru{э́тот результа́т позволя́ет нам} this result enables us to) and \ru{у-до-вле-твор-я́ть} satisfy (\ru{твори́ть} means create); note \ru{вле-} for \ru{вол-} with vowel gradation and metathesis of liquids; see Introduction §5.

\subsection*{7. \ru{да-} --- give}

Compounds of \ru{дава́ть, дать} give (Lat. \emph{dat-}) are:

\begin{center}
\small
\begin{tabular}{@{}>{\cyrillicfont}p{43mm}p{38mm}>{\cyrillicfont}p{57mm}@{}}
задавать $\sim$ задать & prescribe, assign & равно заданному числу\\
задача & problem & фундаментальная задача интегрального исчисления\\
издавать $\sim$ издать & publish & издательство\\
передавать $\sim$ передать & transmit & канал передачи\\
придавать $\sim$ придать & attach, add & придать приращение независимой переменной\\
удаваться $\sim$ удаться & turn out well, succeed & нам удалось построить решение\\
\end{tabular}
\end{center}

\subsection*{8. \ru{дел-} --- divide}

The compounds of \ru{дели́ть $\sim$ раздели́ть} divide conform to the regular pattern. Thus \ru{выделя́ть $\sim$ вы́делить} distinguish, etc.

Sentences illustrating a) \ru{дели́ть $\sim$ раздели́ть} divide (cf. \ru{деле́ние} division), b) \ru{дели́тель} divisor, c) \ru{выдели́ть} distinguish, extract, d) \ru{преде́л} limit, e) \ru{преде́льный} limiting, limit, f) \ru{определя́ть $\sim$ определи́ть} delimit, determine, define (\ru{неопределённый} indefinite), g) \ru{определе́ние} definition, h) \ru{отде́льность} separation (lit. off-divided-ness):

\begin{enumerate}[label=\alph*)]
\item \ru{умно́жим и разде́лим на $(x^2-a^2)$} --- let us multiply and divide by $(x^2-a^2)$;
\item \ru{элеме́нты, для кото́рых $a\ne0$, $b\ne0$, но $ab=0$, называ́ются дели́телями нуля́} --- elements of this kind are called divisors of zero;
\item \ru{мо́жно вы́делить коне́чную подсисте́му} --- it is possible to extract a finite subsystem;
\item \ru{преде́л фу́нкции} --- limit of a function;
\item \ru{то́чка называ́ется преде́льной то́чкой, е́сли в ка́ждой о́крестности \ldots} --- a point is called a limit point if in every neighborhood \ldots;
\item \ru{положе́ние то́чки определя́ется отре́зком} --- the position of the point is determined by the segment;
\item \ru{определе́ние ве́ктора} --- the definition of a vector;
\item \ru{фу́нкция голомо́рфна по ка́ждой переме́нной в отде́льности} --- the function is holomorphic with respect to each variable separately.
\end{enumerate}

\subsection*{9. \ru{зна-} --- know}

In modern Russian the two verbs \ru{знава́ть} and \ru{знать} know are both imperfective, but the compounds of \ru{знава́ть} are imperfective and have the corresponding compounds of \ru{знать} as their perfectives. Thus \ru{признава́ть $\sim$ призна́ть} recognize, with \ru{при́знак} criterion; \ru{узна́ть со́бственное значе́ние} identify an eigenvalue.

The noun \ru{знак} sign, symbol (i.e. means of knowing) gives rise to the verb \ru{зна́чить} (no pf.) to mean with compounds like \ru{че́рез $M_x$ обозна́чим \ldots} by $M_x$ we shall denote \ldots; \ru{что означа́ют э́ти зна́ки?} what do these symbols signify? From \ru{зна́чить} to mean come several important nouns and adjectives, e.g. \ru{значе́ние} meaning, value (e.g. of a variable), \ru{однозна́чный} one-valued, etc.

Sentences illustrating a) \ru{знать} know, b) \ru{знак} sign, c) \ru{зна́ние} knowledge, d) \ru{при́знак} criterion, e) \ru{зна́чить} to mean, f) \ru{значе́ние} value, g) \ru{обозначе́ние} designation, h) \ru{знамена́тель} denominator, i) \ru{однозна́чность} one-valuedness:

\begin{enumerate}[label=\alph*)]
\item \ru{как мы зна́ем} --- as we know;
\item \ru{знак $m$ противополо́жен зна́ку $c$} --- the sign of $m$ is opposite to the sign of $c$;
\item \ru{без зна́ния положе́ния криво́й в простра́нстве} --- without knowledge of the position of the curve in space;
\item \ru{при́знаком перпендикуля́рности ве́кторов явля́ется раве́нство нулю́ их скаля́рного произведе́ния} --- a criterion of perpendicularity is equality of their scalar product to zero;
\item \ru{зна́чит, элеме́нт $a$ явля́ется нулём кольца́} --- it means that the element $a$ is the zero of the ring;
\item \ru{значе́нию $x$ ста́вится в соотве́тствие значе́ние $y$} --- to a value of $x$ there is put in correspondence a value of $y$;
\item \ru{для обозначе́ния произво́дной употребля́ют разли́чные си́мволы} --- various symbols are used for the derivative;
\item \ru{дробь со сравни́тельно небольши́м знамена́телем} --- a fraction with a comparatively small denominator;
\item \ru{отка́жемся от тре́бования однозна́чности} --- we discard the requirement of one-valuedness.
\end{enumerate}

\subsection*{10. \ru{зыв-} --- call}

The simple verb is \ru{звать $\sim$ позва́ть} call. The lexical compounds commonest in mathematics are \ru{вызыва́ть $\sim$ вы́звать} call up, evoke, give rise to (e.g. \ru{вы́звать оши́бку} cause an error) and \ru{называ́ть $\sim$ назва́ть} call, e.g. by definition.

Sentences illustrating a) \ru{вызыва́ть} cause, b) \ru{называ́ть} call, c) \ru{назва́ние} name:

\begin{enumerate}[label=\alph*)]
\item \ru{ма́лая оши́бка в да́нных мо́жет вызыва́ть большу́ю оши́бку} --- a small error in the data can cause a large error;
\item \ru{ра́зность расстоя́ний от двух то́чек, называ́емых фо́кусами} --- the difference of the distances from two points called foci;
\item \ru{перечи́слим не́которые кла́ссы фу́нкций, получи́вших назва́ние элемента́рных} --- let us enumerate some classes of functions having received the name ``elementary''.
\end{enumerate}

\subsection*{11. \ru{иск-} --- seek}

The compounds from \ru{иска́ть} seek (cognate with the English \emph{ask}) fit precisely into the scheme of Chapter IV; e.g. (for \ru{-ыск-} in place of \ru{-иск-} see Chap. II §8): \ru{оты́скивать $\sim$ отыска́ть} search out; \ru{разы́скивать $\sim$ разыска́ть} search in all directions, investigate, as in \ru{разыска́ние фу́нкций таки́х, что́бы \ldots} the search for functions such that \ldots

Sentences illustrating a) \ru{иска́ть} seek, b) \ru{отыскание} search:

\begin{enumerate}[label=\alph*)]
\item \ru{э́то и есть иско́мое уравне́ние} --- this is the sought-for equation;
\item \ru{Евкли́ду принадлежи́т ме́тод отыска́ния наибо́льшего о́бщего дели́теля} --- Euclid is credited with a method for finding the greatest common divisor.
\end{enumerate}

\subsection*{12. \ru{каз-} --- show}

The simple verb \ru{*каза́ть} show is obsolete, but its reflexive \ru{каза́ться} (show itself, appear to be) and many of its compounds are very common, e.g. \ru{вы́сказывать $\sim$ вы́сказать} speak out, assert, state (\ru{исчисле́ние выска́зываний} calculus of statements); \ru{дока́зывать $\sim$ доказа́ть} prove (\ru{доказа́тельство} proof; \ru{что и тре́бовалось доказа́ть} which in fact it was required to prove); \ru{ока́зываться $\sim$ оказа́ться} turn out (\ru{ока́зывается, что зада́ча неразреши́ма}); \ru{отка́зываться $\sim$ отказа́ться} deny oneself, relinquish (\ru{отказа́ться от тре́бования}); \ru{пока́зывать $\sim$ показа́ть} expound, show (\ru{Евкли́д показа́л, что \ldots}; \ru{показа́тельная фу́нкция} exponential function); \ru{(говори́ть) $\sim$ сказа́ть} say (\ru{так сказа́ть}); \ru{ска́зываться} express itself (\ru{оши́бка ска́зывается}); \ru{ука́зывать $\sim$ указа́ть} indicate (\ru{число́ ука́занного ви́да}).

\subsection*{13. \ru{кры-} --- cover}

Here the simple verb \ru{крыть} to cover remains in use, and its (still common) perfective \ru{покры́ть} has formed a new imperfective \ru{покрыва́ть}. Its commonest mathematical compound is \ru{открыва́ть $\sim$ откры́ть} uncover, open; compare \ru{откры́тое покры́тие} an open covering and \ru{е́сли за́мкнутый промежу́ток покрыва́ется бесконе́чной систе́мой откры́тых промежу́тков} if a closed interval is covered by an infinite system of open intervals.\corrnote{原文称 the simple verb \ru{крыть} to cover is obsolete。该判断不正确：\ru{крыть} 仍是标准俄语动词。此处只改正 obsolete 这一点，其余构词说明按原文保留。}

\subsection*{14. \ru{мен-} --- change}

The simple verb \ru{меня́ть} change (with perfectives \ru{обменя́ть} and \ru{поменя́ть}) has various (perfective) lexical compounds \ru{выменя́ть} exchange, \ru{разменя́ть} disperse, etc. that form their imperfectives \ru{выме́нивать, разме́нивать}, etc. regularly by the first method, i.e. with the suffix \ru{-ива-} (see Chap. IV §4). But most of the mathematically important compounds, e.g. \ru{заменя́ть} substitute, \ru{переменя́ть} vary, \ru{применя́ть} apply, have come to be regarded as imperfectives (and have been provided with perfectives \ru{замени́ть}, etc.) because they happen to have the same ending \ru{-ять} as is regularly used to imperfectivize verbs in \ru{-ить} (again see Chap. IV §4). Thus \ru{заменя́ть $\sim$ замени́ть} replace, substitute; \ru{изменя́ть $\sim$ измени́ть} alter; \ru{переменя́ть $\sim$ перемени́ть} vary; \ru{применя́ть $\sim$ примени́ть} adapt, apply.

Sentences illustrating a) \ru{замени́ть} replace, b) \ru{измене́ние} variation, c) \ru{переме́нная} variable, d) \ru{примени́ть} apply:

\begin{enumerate}[label=\alph*)]
\item \ru{в результа́те теоре́мы Туэ--Зи́геля--Ро́та нельзя́ замени́ть $1/b^{2+\varepsilon}$ че́рез $1/b^2$} --- the exponent cannot be sharpened by this replacement;
\item \ru{фу́нкция с ограни́ченным измене́нием} --- a function of bounded variation;
\item \ru{пусть даны́ две переме́нные $x$ и $y$} --- let there be given two variables $x$ and $y$;
\item \ru{приме́ним пра́вило подстано́вки} --- we shall apply the rule of substitution.
\end{enumerate}

\subsection*{15. \ru{мер-} --- measure}

Here the verb \ru{ме́рить} measure and its compounds fit precisely into the scheme of Chapter IV; e.g. \ru{ме́рить $\sim$ поме́рить} to measure, \ru{ме́ра} measure, \ru{по ме́ньшей ме́ре} at least (lit. by least measure); \ru{измеря́ть $\sim$ измери́ть} to measure, \ru{несоизмери́мость} incommensurability; \ru{примеря́ть $\sim$ приме́рить} to fit, try on, \ru{приме́р} example, \ru{наприме́р} for example.

Sentences illustrating a) \ru{измеря́ть} measure, b) \ru{измере́ние} dimension, c) \ru{трёхме́рный} three-dimensional (cf. \ru{дву́хмерный} two-dimensional), d) \ru{многоме́рный} n-dimensional, e) \ru{разме́рность} dimensionality:

\begin{enumerate}[label=\alph*)]
\item \ru{зная пе́рвую квадрати́чную фо́рму, мо́жно измеря́ть дли́ны, углы́ и пло́щади на пове́рхности};
\item \ru{число́ измере́ний ме́ньше четырёх};
\item \ru{мно́жество ве́кторов трёхме́рного евкли́дова простра́нства};
\item \ru{ри́манова геоме́трия --- многоме́рное обобще́ние геоме́трии на пове́рхности};
\item \ru{разме́рность э́того простра́нства называ́ется кра́тностью со́бственного значе́ния}.
\end{enumerate}

\subsection*{16. \ru{мог-} --- be able}

\ru{мог-} be able (cf. \emph{may, might}). The simple verb \ru{мочь $\sim$ смочь} be able forms the lexical compound \ru{помога́ть $\sim$ помо́чь} to help (note this rare use of \ru{по-} in an aspectual compound, i.e. other than as a mere perfectivizer).

Sentences illustrating a) \ru{мочь} can, b) \ru{мо́жно} possible, c) \ru{возмо́жный} possible, \ru{возмо́жно} possibly, as ... as possible, \ru{возмо́жность} possibility, d) \ru{мо́щность} power, e) \ru{невозмо́жно} impossible, f) \ru{вспомога́тельный} auxiliary, g) \ru{по́мощь} help:

\begin{enumerate}[label=\alph*)]
\item \ru{кольцо́ с дели́телями нуля́ не мо́жет содержа́ться в по́ле};
\item \ru{мо́жно счита́ть э́ти чи́сла ко́синусом и си́нусом не́которого угла́};
\item \ru{выбира́ют рациона́льное число́ возмо́жно про́стое} --- they choose a rational number as simple as possible;
\item \ru{об эквивале́нтных мно́жествах говоря́т, что они́ име́ют одина́ковую мо́щность};
\item \ru{что невозмо́жно};
\item \ru{рассмо́трим вспомога́тельную пряму́ю};
\item \ru{разры́вная фу́нкция Дирихле́ не мо́жет быть изображена́ при по́мощи одного́ перехо́да к преде́лу}.
\end{enumerate}

\subsection*{17. \ru{рез-} --- cut}

Both \ru{рез-} and \ru{сек-} (No. 18 following) mean to cut, \ru{рез-} to cut around and \ru{сек-} to cut through. The verb \ru{ре́зать} cut and its compounds fit precisely into our pattern; thus \ru{ре́зать $\sim$ заре́зать} cut; \ru{реша́ть $\sim$ реши́ть} decide, solve (note the consonant variation), with \ru{резе́ц} chisel and \ru{реше́ние} decision, solution; \ru{выреза́ть $\sim$ вы́резать} cut out; \ru{обреза́ть $\sim$ обре́зать} cut around, sculpt; \ru{отреза́ть $\sim$ отре́зать} cut off; \ru{разреза́ть $\sim$ разре́зать} to slit.

The vowel-grade \ru{-а-} (as above in \ru{образ}) also occurs in such verbs as \ru{рази́ть $\sim$ порази́ть} cut down, strike (the derived noun \ru{раз} a time means that which is struck or marked off); \ru{выража́ть $\sim$ вы́разить} formulate, sketch out, express; \ru{отража́ть $\sim$ отрази́ть} reflect. Compare \ru{вы́резка} engraving, \ru{о́браз} (sculpted) image, form, manner, \ru{отре́зок прямо́й} segment of a line, \ru{разре́з в пове́рхности} a slit in the plane, \ru{фо́рмула выража́ет теоре́му}, \ru{тео́рия отраже́ния}.

From the noun \ru{о́браз} image, form, manner, comes the verb \ru{*образи́ть} to form (obsolete as an uncompounded verb) with the important (regularly constructed) compounds \ru{изобража́ть $\sim$ изобрази́ть} represent (\ru{изображе́ние по Лапла́су} Laplace transform) and \ru{отобража́ть $\sim$ отобрази́ть} to map (\ru{конфо́рмное отображе́ние} conformal mapping). The simple verb \ru{*образи́ть} has been replaced by \ru{образо́вывать $\sim$ образова́ть} form, generate (\ru{образу́ющая ко́нуса} generator of a cone), which is constructed from the noun \ru{о́браз} with the verbal suffix \ru{-овать} and forms the further compound \ru{преобразо́вывать $\sim$ преобразова́ть} to transform (\ru{преобразова́ние Лапла́са} Laplace transformation).

Sentences illustrating a) \ru{о́браз} image, form, manner, b) \ru{изобрази́ть} represent, illustrate, c) \ru{вы́разить} express, d) \ru{выраже́ние} expression, e) \ru{отображе́ние} mapping, f) \ru{образова́ть} to form, g) \ru{проо́браз} preimage, h) \ru{многообра́зие} manifold:

\begin{enumerate}[label=\alph*)]
\item \ru{аналоги́чным о́бразом; таки́м о́бразом};
\item \ru{построе́ние су́ммы ве́кторов изображено́ на рису́нке};
\item \ru{на терминоло́гии после́довательностей непреры́вность вы́разится так:};
\item \ru{выраже́ние $A\cdot dx$ называ́ется дифференциа́лом};
\item \ru{взаи́мно однозна́чные отображе́ния};
\item \ru{па́ры це́лых чи́сел образу́ют кольцо́, е́сли \ldots};
\item \ru{вся́кий элеме́нт облада́ет, вообще́ говоря́, мно́гими разли́чными проо́бразами};
\item \ru{не́которое многообра́зие \ldots}
\end{enumerate}

\subsection*{18. \ru{сек-} --- cut}

The verb \ru{сечь} to cut is similar in form to \ru{мочь} to be able. Thus \ru{сечь $\sim$ вы́сечь} cut (cf. \emph{secant, section}); \ru{секу́щая} (pres. participle) cutting, secant; \ru{отсека́ть} cut off; \ru{пересека́ть $\sim$ пересе́чь} intersect. Compare \ru{поня́тие о сече́нии}, \ru{преде́льное положе́ние секу́щей}, \ru{пряма́я отсека́ет отре́зки на ося́х}, \ru{пряма́я пересека́ет ось $Oy$}, \ru{крива́я име́ет самопересече́ния}.

Sentences illustrating a) \ru{сечь} to cut, b) \ru{сече́ние} a cut, c) \ru{отсека́ть} cut off, d) \ru{пересека́ть} intersect, e) \ru{пересече́ние} intersection:

\begin{enumerate}[label=\alph*)]
\item \ru{преде́льное положе́ние секу́щей называ́ют каса́тельной};
\item \ru{тео́рия Дедеки́нда о сече́ниях};
\item \ru{отре́зок, отсека́емый на о́си};
\item \ru{разбие́ние мно́жества на непересека́ющиеся подмно́жества};
\item \ru{е́сли крива́я име́ет самопересече́ние}.
\end{enumerate}

\subsection*{19. \ru{след-} --- slide, leave a footprint, trace, follow}

\ru{следи́ть} (no pf.) follow, keep watch on; \ru{след} trace, as in \ru{след ма́трицы} trace of a matrix; \ru{сле́дующий} the following; \ru{иссле́довать} (its own pf.) investigate, as in \ru{иссле́дуем фо́рму э́ллипса} let us investigate the form of an ellipse.

Various derived words are \ru{сле́довательно} consequently, \ru{сле́дствие} corollary, \ru{после́довательность} sequence, \ru{после́довательный} successive, \ru{после́довательно} successively, \ru{после́дний} last, \ru{всле́дствие} as a consequence of, as in the following examples.

Sentences illustrating a) \ru{сле́довать} follow, b) \ru{сле́дствие} corollary, c) \ru{после́довательность} sequence, d) \ru{после́дний} last, e) \ru{всле́дствие} as a consequence of, f) \ru{после́довательно} successively:

\begin{enumerate}[label=\alph*)]
\item \ru{из определе́ния сле́дует, что \ldots};
\item \ru{име́ем сле́дствие};
\item \ru{на терминоло́гии после́довательностей непреры́вность вы́разится так:};
\item \ru{интегра́л от после́днего выраже́ния вычисля́ется так:};
\item \ru{всле́дствие непреры́вности};
\item \ru{мы после́довательно определи́м це́лые, рациона́льные, де́йствительные и ко́мплексные чи́сла}.
\end{enumerate}

\subsection*{20. \ru{чит-} --- consider}

Some words derived from this root still have the general meaning consider, reckon, but two special meanings have developed: read and calculate. Thus \ru{счита́ть} consider, reckon as; \ru{чита́ть $\sim$ прочита́ть} read; \ru{рассчи́тывать $\sim$ рассчита́ть} calculate; \ru{вычита́ть $\sim$ вы́честь} subtract; \ru{число́} (orig. \ru{чит-сло}) number; \ru{числово́й} numerical; \ru{числи́тель} numerator; \ru{вычисля́ть $\sim$ вы́числить} compute; \ru{перечисля́ть $\sim$ перечисли́ть} count through, enumerate; \ru{исчисле́ние} calculus; \ru{вы́чет} residue; \ru{учёт} account; \ru{счётный} countable.

Sentences illustrating a) \ru{счита́ть} consider, b) \ru{вы́чет} residue, c) \ru{счётный} and \ru{несчётный} countable and uncountable, d) \ru{вычита́ние} subtraction, e) \ru{исчисле́ние} calculus, f) \ru{рассчита́ть} calculate:

\begin{enumerate}[label=\alph*)]
\item \ru{обозна́чим че́рез $(\rho,\theta)$ поля́рные координа́ты то́чки $M(x,y)$, счита́я поля́рной о́сью $Ox$};
\item \ru{приме́ром по́ля характери́стики $p$ явля́ется кольцо́ вы́четов по мо́дулю $p$};
\item \ru{мно́жество, не явля́ющееся ни коне́чным, ни счётным, называ́ется несчётным};
\item \ru{вычита́нием натура́льных чи́сел называ́ется опера́ция, обра́тная сложе́нию};
\item \ru{теоре́ма о сре́днем значе́нии в дифференциа́льном исчисле́нии};
\item \ru{рассчита́ем вре́мя, необходи́мое для того́, что \ldots}
\end{enumerate}

\section{Forty-one roots of moderate productivity}

Each of the following 41 roots produces several words of mathematical importance; here we give at least one pair of aspect-partners for each root.\corrnote{原文标题与首句均写 ``40 roots''；实际逐项有 41 个根词标题，故仅改正计数。}

\begingroup
\small
\setlength{\tabcolsep}{3pt}
\renewcommand{\arraystretch}{1.08}
\begin{longtable}{@{}>{\cyrillicfont}p{23mm}p{28mm}p{92mm}@{}}
\toprule
\normalfont Root & Meaning & Representative derivatives and mathematical uses\\
\midrule
\endhead
би- & beat, strike & \ru{разбива́ть $\sim$ разби́ть} decompose/partition; \ru{разбие́ние} decomposition\\
велик- & big & \ru{увели́чивать $\sim$ увели́чить} increase; \ru{величина́} magnitude\\
верх- & top & \ru{соверша́ть $\sim$ соверши́ть} complete; \ru{соверше́нно}, \ru{верши́на}, \ru{ве́рхний}, \ru{пове́рхность}, \ru{све́рху}\\
влад- & possess & \ru{облада́ть} possess; \ru{о́бласть} domain\\
вяз- & knot, bind & \ru{вяза́ть $\sim$ связа́ть}; \ru{обя́зывать $\sim$ обяза́ть}; \ru{обяза́тельно}; \ru{свя́зка}; \ru{связь}\\
гран- & face, border & \ru{ограни́чивать $\sim$ ограни́чить}; \ru{грани́ца}; \ru{грани́чный}; \ru{ограниче́ние}; \ru{ограни́ченность}; \ru{неограни́ченный}\\
грех- & mistake & \ru{погреша́ть $\sim$ погреши́ть}; \ru{погре́шность} error\\
дел- & do & \ru{де́лать $\sim$ сде́лать}; \ru{де́ло}; distinct from \ru{дел-} ``divide''\\
держ- & hold & \ru{содержа́ть}; \ru{содержа́ние} content\\
доб- & suitable & \ru{понадо́биться} be necessary; \ru{подо́бие}; \ru{подо́бный}; \ru{удо́бство}; \ru{неудо́бный}\\
кас- & touch & \ru{соприкаса́ться $\sim$ соприкосну́ться}; \ru{каса́тельная}; \ru{каса́ние}; \ru{соприкаса́ющаяся пло́скость}\\
клон- & bend, incline & \ru{клони́ть}; \ru{отклоня́ть $\sim$ отклони́ть}; \ru{наклоня́ть $\sim$ наклони́ть}; \ru{накло́н}\\
ключ- & close & \ru{включа́ть $\sim$ включи́ть}; \ru{включе́ние}; \ru{включи́тельно}; \ru{заключа́ть $\sim$ заключи́ть}; \ru{заключе́ние}; \ru{исключа́ть $\sim$ исключи́ть}\\
кон- & end & \ru{конча́ть $\sim$ ко́нчить}; \ru{коне́ц}; \ru{коне́чный}; \ru{бесконе́чный}; \ru{коне́чность}; \ru{оконча́тельный}; \ru{наконе́ц}\\
круг- & circle & \ru{окружа́ть $\sim$ окружи́ть}; \ru{окру́жность}; \ru{круг}\\
крут- & twist & \ru{крути́ть $\sim$ закрути́ть}; \ru{круче́ние} torsion\\
лег-/льз- & ease, use & \ru{по́льзоваться $\sim$ воспо́льзоваться}; \ru{испо́льзовать}; \ru{поле́зный}; \ru{нельзя́}; \ru{лёгкий}; \ru{легко́}\\
лик- & face & \ru{отлича́ть $\sim$ отличи́ть}; \ru{отли́чие}; \ru{отли́чный}; \ru{разли́чный}; \ru{разли́чие}\\
меш- & mix & \ru{сме́шивать $\sim$ смеша́ть}; \ru{сме́шанная гру́ппа} mixed group\\
мест- & place & \ru{перемеща́ть $\sim$ перемести́ть}; \ru{ме́сто}; \ru{вме́сте}; \ru{вме́сто}; \ru{совме́стный} consistent/compatible\\
мык- & close & \ru{замыка́ть $\sim$ замкну́ть}; \ru{замыка́ние}; \ru{за́мкнутый}; \ru{примыка́ть $\sim$ примкну́ть}; \ru{примыка́ние}\\
нов- & new & \ru{обновля́ть $\sim$ обнови́ть}; \ru{но́вый}\\
пад- & fall & \ru{па́дать $\sim$ упа́сть}; \ru{совпада́ть $\sim$ совпа́сть} coincide\\
плоск- & flat & \ru{уплоща́ть $\sim$ уплости́ть}; \ru{пло́скость}; \ru{пло́ский}; \ru{пло́щадь}\\
плот- & thick, dense & \ru{уплотня́ть $\sim$ уплотни́ть}; \ru{пло́тный} dense\\
прос- & ask & \ru{проси́ть $\sim$ попроси́ть}; \ru{вопро́с} question\\
прост- & simple & \ru{упроща́ть $\sim$ упрости́ть}; \ru{просто́й} simple, prime\\
пряг- & harness, join & \ru{сопряга́ть $\sim$ сопря́чь}; \ru{сопряжённый опера́тор}; \ru{самосопряжённый}\\
прям- & straight & \ru{спрямля́ть $\sim$ спрями́ть}; \ru{спрямля́ющая пло́скость}; \ru{пряма́я}\\
раст- & grow & \ru{прираста́ть $\sim$ прирасти́}; \ru{прираще́ние}; \ru{рост}; \ru{возраста́ть $\sim$ возрасти́}; \ru{возраста́ние}\\
реч- & speech & \ru{противоре́чить}; \ru{противоре́чие}; \ru{отрица́тельность} negativity\\
род- & birth, kind & \ru{порожда́ть $\sim$ породи́ть}; \ru{род}; \ru{одноро́дный}; \ru{неодноро́дный}\\
рыв- & break, tear & \ru{рвать}; \ru{разры́в}; \ru{разры́вный}; \ru{непреры́вный}; \ru{непреры́вно}; \ru{непреры́вность}\\
ряд- & row, series & \ru{упоря́дочивать $\sim$ упоря́дочить}; \ru{поря́док}\\
смотр- & look & \ru{смотре́ть $\sim$ посмотре́ть}; \ru{рассма́тривать $\sim$ рассмотреть}\\
стро- & construct & \ru{стро́ить $\sim$ постро́ить}; \ru{строе́ние}; \ru{построе́ние}\\
треб- & demand & \ru{тре́бовать $\sim$ потребова́ть}; \ru{тре́бование}; \ru{употребля́ть $\sim$ употреби́ть} use\\
тык- & pierce, point & \ru{ты́кать $\sim$ ткну́ть}; \ru{то́чка}; \ru{то́чность}; \ru{то́чный}; \ru{то́чно}; \ru{уточня́ть $\sim$ уточни́ть}; \ru{то́чечно} pointwise\\
ук- & learn & \ru{изуча́ть $\sim$ изучи́ть}; \ru{нау́ка}; \ru{обы́чный}; \ru{обы́чно}; \ru{обыкнове́нный}\\
ча- & begin & \ru{начина́ть $\sim$ нача́ть}; \ru{нача́ло} beginning/origin; \ru{нача́льный} initial\\
част- & part & \ru{уча́ствовать}; \ru{часть}; \ru{уча́сток}; \ru{ча́стность}; \ru{ча́стный} partial\\
\bottomrule
\end{longtable}
\endgroup

\section{Other nouns and adjectives}

Finally, here are a few nouns and adjectives which, though most of them are fairly important in mathematics, are not readily associated with verbs according to the system of Chapters IV and V.

\subsection*{Masculine nouns}
\noindent\ru{ко́рень} root; \ru{рису́нок} diagram; \ru{у́гол} angle; \ru{шар} sphere, ball.

\subsection*{Feminine nouns}
\noindent\ru{длина́} length; \ru{дробь} fraction; \ru{кра́тность} multiplicity; \ru{крива́я} curve; \ru{кривизна́} curvature; \ru{о́крестность} neighbourhood; \ru{ось} axis; \ru{оши́бка} error; \ru{ско́бка} bracket; \ru{ско́рость} speed; \ru{сте́пень} degree; \ru{це́лостность} integrity.

\subsection*{Neuter nouns}
\noindent\ru{кольцо́} ring (\ru{подкольцо́} subring); \ru{по́ле} field (\ru{подпо́ле} subfield); \ru{простра́нство} space (\ru{простра́нственный} space adj.; \ru{подпростра́нство} subspace); \ru{сло́во} word, with the derivative \ru{усло́вие} condition.

\subsection*{Adjectives}
\noindent\ru{гла́вный} principal; \ru{действи́тельный} real (\ru{действи́тельно} in fact); \ru{дро́бный} fractional; \ru{друго́й} other; \ru{любо́й} arbitrary; \ru{ма́лый} small (\ru{ма́ло} by a small amount); \ru{неизве́стный} unknown; \ru{о́бщий} general; \ru{сла́бый} weak (\ru{сла́бо} weakly); \ru{ста́рый} old; \ru{углово́й} angular; \ru{це́лый} integral, whole; \ru{цепно́й} chain-like, continued (from \ru{цепь} chain).

\fussy
